\section{Discussion}
\label{sec:discussion}

% Extended discussion content (CFL echo-back, rate-distortion analysis,
% LLMLingua comparison detail, NoTool bias, latency analysis, scope table,
% extended threats to validity, future directions) cut from main text.

\subsection{Three Mechanisms of TSCG Improvement}
\label{sec:three-mechanisms}

TSCG's accuracy improvements arise from three distinct mechanisms.
\textbf{M1: Format Translation} converts JSON schemas to structured text---the dominant mechanism for Class~1 models, where JSON parsing is the bottleneck (e.g., Phi-4: 0\%$\to$90\%, entirely format-driven per E1/E4).
\textbf{M2: Structural Reorganization} (CAS reordering, CFL constraint positioning, CFO causal ordering) improves accuracy beyond format change---the dominant mechanism for Class~2 frontier models ($+$10.9\,pp on Claude Sonnet~4 in text-mode, $+$24.4\,pp on GPT-5.2).
\textbf{M3: Token Reduction} decreases attention dilution as a secondary effect; GPT-4o achieves positive accuracy improvement even with negative token savings on certain scenarios, proving M2 operates independently of M3.

The mechanisms contribute differently per model class: M1 dominates for small models (Class~1), M2 for frontier models (Class~2), and neither contributes meaningfully for neutral models (Class~3). For Qwen architectures (Class~4), balanced M2 transforms actively harm accuracy, but conservative filler removal preserves benefits.

\subsection{Four-Class Taxonomy}
\label{sec:enabler-optimizer}

E4 text-baseline experiments (2{,}520 calls, 6~small models), N3 conservative ablation (180 calls), and N1 30B benchmarks (840 calls, 2~models) yield a four-class taxonomy across 12~models (Table~\ref{tab:four-class}):

\begin{table}[h]
\centering
\small
\caption{Four-class taxonomy of TSCG effects (12 models, 4B--32B + frontier). Text~$\Delta$: compression gain (TSCG minus text baseline).}
\label{tab:four-class}
\begin{tabular}{llrrl}
\toprule
\textbf{Class} & \textbf{Models} & \textbf{JSON $\Delta$} & \textbf{Text $\Delta$} & \textbf{Recommendation} \\
\midrule
1: Format-dom. & Phi-4, Mistral 7B, Gemma 4B, Qwen3 4B & +17 to +90\,pp & $-$7 to $-$23\,pp & Conservative \\
2: Compression & Claude, GPT-4o, GPT-5.2 & +5 to +11\,pp & +5 to +11\,pp & Balanced \\
3: Neutral & Llama 8B, Gemma 12B, Mistral-Sm.\ 24B & +5 to +9\,pp & $\approx$0\,pp & Conservative \\
4: Cons.-only & Qwen3 14B, Qwen2.5-Coder 32B & $-$5 to $-$7\,pp & $-$5 to $-$13\,pp & Conservative only \\
\bottomrule
\end{tabular}
\end{table}

\textbf{Class~1} models gain dramatically from JSON-baseline TSCG but show \emph{negative} compression against text baselines---the benefit is format translation (JSON$\to$text), confirmed by E4 decomposition.
\textbf{Class~2} frontier models show genuine structural compression ($+$5--11\,pp) that persists when format confound is eliminated.
\textbf{Class~3} models are robust across all formats; TSCG neither helps nor harms.
\textbf{Class~4} Qwen models degrade under balanced TSCG but improve $+$4.4\,pp under conservative SDM (N3)---the degradation is architecture-specific, persisting from 14B to 32B (N1), implicating CAS/DRO structural transforms rather than capacity limitations.

\paragraph{Deployment guidance.} Since every production API transmits JSON: (1)~Class~1 models achieve 0--49\% at $>$15 tools but 65--90\% under TSCG---format translation is the needed intervention; (2)~Class~2 benefits from balanced profiles; (3)~Class~4 requires conservative mode exclusively; (4)~Class~3 can use TSCG without risk.

\subsection{Naive Truncation Honesty}
\label{sec:naive-truncation-discussion}

At 16~well-named tools (Scenario~A), naive truncation (retaining only \texttt{ToolName(param:type,...)}) matches TSCG: 87.0\% vs.\ 87.8\% overall.
Claude Sonnet~4 recognizes tools by distinctive names, making descriptions dispensable.
The gap emerges at scale: at 50~tools with ambiguous names (Scenario~C-50), TSCG achieves 100\% while naive truncation drops to 98.5\%, concentrated in multi-tool sequencing (87.5\% vs.\ 75.0\%, $-$12.5\,pp) and parameter extraction (100\% vs.\ 92.5\%).
TSCG's structural reorganization preserves the semantic cues needed for disambiguation---parameter descriptions, usage constraints, return types---that naive truncation discards.

\subsection{Infrastructure Impact}
\label{sec:infrastructure-discussion}

TSCG reduces per-call schema overhead from $O(n \cdot k)$ to $O(n \cdot k/c)$ where $c \approx 3.5\times$ for structured schemas.
At production volumes (100k calls/day with 16~tools at \$3/MTok), the savings exceed \$30{,}000/month.
Combined with sub-millisecond deterministic execution, zero dependencies, and a 27.7\,KB bundle, TSCG deploys as transparent middleware in any agent framework---Claude Code, LangChain, CrewAI, or MCP servers---without GPU infrastructure or external API calls.

For small models (4B--14B), the infrastructure impact extends beyond cost: TSCG enables deployment of local models as functional tool-use agents in edge, cost-sensitive, and privacy-constrained environments where frontier APIs are unavailable or unacceptable.
